% bibliography of Volume 2: the 326 of 590 entries of ch/bib.tex it cites; generated
\begin{thebibliography}{999}\small\setlength{\itemsep}{0pt}\setlength{\parsep}{0pt}\setlength{\parskip}{0pt}
% verified: CrossRef 10.1002/nme.995 2026-10-09

% verified: CrossRef 10.1007/978-1-4757-8137-3_5 2026-10-24
\bibitem{ABR01} S. Axler, P. Bourdon and W. Ramey. \emph{Harmonic Function Theory}, second edition. Springer, 2001.
% verified: CrossRef 10.1145/103418.103439 2026-10-24
\bibitem{AK91} D. Applegate and R. Kannan. Sampling and integration of near log-concave functions. In \emph{Proceedings of the 23rd Annual ACM Symposium on Theory of Computing}, 1991, 156--163.
% verified: pending (not found: CrossRef and OpenAlex searched by title; unpublished technical report)
\bibitem{AKP90} D. Applegate, R. Kannan and N. G. Polson. Random polynomial time algorithms for sampling from joint distributions. Technical report, Department of Statistics, Carnegie Mellon University, 1990.
% verified: CrossRef 10.1111/j.2517-6161.1974.tb00989.x 2026-10-14; the record gives the later journal title, Series B: Statistical Methodology
\bibitem{AndrewsMallows74} D. F. Andrews and C. L. Mallows. Scale mixtures of normal distributions. \emph{Journal of the Royal Statistical Society Series B (Methodological)} 36 (1974), 99--102.
% verified: zbMATH 0456.60081 2026-09-27 (CrossRef 10.1137/1125045 dates the translated issue 25(2) January 1981; Russian original Teor. Veroyatnost. i Primenen. 25 (1980), 366)
\bibitem{Anulova80} S. V. Anulova. On Markov stopping times with a given distribution for a Wiener process. \emph{Theory of Probability and its Applications} 25 (1980), 362--366.
% verified: CrossRef 10.1090/s0025-5718-2010-02426-3 2026-09-27 (CrossRef splits the surname as 'de Reyna' and dates it online 2010; the AMS page and zbMATH 1228.11127 give Arias de Reyna, 80 (2011))
\bibitem{AriasDeReyna11} J. Arias de Reyna. High precision computation of Riemann's zeta function by the Riemann-Siegel formula, I. \emph{Mathematics of Computation} 80 (2011), 995--1009.
% verified: CrossRef 10.1007/bf02786970 2026-10-23
\bibitem{ASW52} M. Aissen, I. J. Schoenberg and A. M. Whitney. On the generating functions of totally positive sequences I. \emph{Journal d'Analyse Math\'ematique} 2 (1952), 93--103.
% verified: zbMATH 45.0719.02 2026-09-27 (JFM title reads 'Sur las'; title and pages 1979 to 1981 checked on the scan of C. R. 158, seance du 29 juin 1914)
\bibitem{Backlund14} R. J. Backlund. Sur les z\'eros de la fonction $\zeta(s)$ de Riemann. \emph{Comptes Rendus de l'Acad\'emie des Sciences} 158 (1914), 1979--1981.
% verified: CrossRef 10.1007/bf02422950 2026-09-26
\bibitem{Backlund} R. J. Backlund. \"Uber die Nullstellen der Riemannschen Zetafunktion. \emph{Acta Mathematica} 41 (1918), 345--375.
% verified: CrossRef 10.1214/aop/1176996548 2026-10-23 (Ann. Probab. 2, 792-804; OpenAlex agrees)
\bibitem{BalkemaDeHaan74} A. A. Balkema and L. de Haan. Residual life time at great age. \emph{Annals of Probability} 2 (1974), 792--804.
% verified: CrossRef 10.1016/0022-2364(87)90023-0 2026-10-24
\bibitem{Barkhuijsen87} H. Barkhuijsen, R. de Beer and D. van Ormondt. Improved algorithm for noniterative time domain model fitting to exponentially damped magnetic resonance signals. \emph{Journal of Magnetic Resonance} 73 (1987), 553--557.
% verified: pending (as above)
\bibitem{BarlowProschan75} R. E. Barlow and F. Proschan. \emph{Statistical Theory of Reliability and Life Testing}. Holt, Rinehart and Winston, 1975.
% verified: CrossRef 10.1016/s0377-0427(00)00336-8 2026-10-25 (OpenAlex agrees on DOI and pages)
\bibitem{BBC00} J. M. Borwein, D. M. Bradley and R. E. Crandall. Computational strategies for the Riemann zeta function. \emph{Journal of Computational and Applied Mathematics} 121 (2000), no. 1--2, 247--296. doi:10.1016/s0377-0427(00)00336-8.
% verified: CrossRef 10.1287/moor.2016.0817 2026-10-24
\bibitem{BBT17} H. H. Bauschke, J. Bolte and M. Teboulle. A descent lemma beyond Lipschitz gradient continuity: first-order methods revisited and applications. \emph{Mathematics of Operations Research} 42 (2017), 330--348.
% verified: CrossRef 10.1214/15-AAP1113 2026-10-25
\bibitem{BJKT16} A. Beskos, A. Jasra, N. Kantas and A. Thiery. On the convergence of adaptive sequential Monte Carlo methods. \emph{Annals of Applied Probability} 26 (2016), 1111--1146.
% verified: arXiv 1306.0856 abstract page 2026-10-25 (journal, volume, issue, pages and DOI as recorded there)
\bibitem{BLM14} H. M. Bui, S. J. Lester and M. B. Milinovich. On Balazard, Saias, and Yor's equivalence to the Riemann Hypothesis. \emph{Journal of Mathematical Analysis and Applications} 409 (2014), no. 1, 244--253. doi:10.1016/j.jmaa.2013.06.063, arXiv:1306.0856.
% verified: CrossRef 10.1090/s0025-5718-07-01999-0 2026-10-11 (Math. Comp. 76, no. 260, 2045-2049); abstract (CrossRef and OpenAlex): zeros off the critical line of the Dirichlet series of Davenport and Heilbronn
\bibitem{BSO07} E. Balanzario and J. S\'anchez-Ortiz. Zeros of the Davenport--Heilbronn counterexample. \emph{Mathematics of Computation} 76 (2007), 2045--2049.
\bibitem{BNKS82} O. E. Barndorff-Nielsen, J. Kent and M. S{\o}rensen. Normal variance-mean mixtures and $z$ distributions. \emph{International Statistical Review} 50 (1982), 145 ff.
% verified: CrossRef 10.2977/prims/31 2026-10-09 (the CrossRef abstract states the Sato Tate conjecture for non-CM holomorphic cuspidal newforms of weight greater than one)
\bibitem{BLGHT11} T. Barnet-Lamb, D. Geraghty, M. Harris and R. Taylor. A family of Calabi--Yau varieties and potential automorphy II. \emph{Publications of the Research Institute for Mathematical Sciences} 47 (2011), 29--98.
% verified: CrossRef 10.1090/mcom/3599 2026-10-10 (title, journal, volume 90, pages 1455 to 1482, issued 2021, authors Bennett, Martin, O'Bryant, Rechnitzer); arXiv 2005.02989 abstract checked for the bound used in Theorem thm:dbeta:minall(v): for conductor q > 1, T >= 5/7 and l = log(q(T+2)/(2 pi)) > 1.567, |N(T,chi) - ((T/pi) log(qT/(2 pi e)) - chi(-1)/4)| <= 0.22737 l + 2 log(1+l) - 0.5
\bibitem{BMOR21} M. A. Bennett, G. Martin, K. O'Bryant and A. Rechnitzer. Counting zeros of Dirichlet $L$-functions. \emph{Mathematics of Computation} 90 (2021), 1455--1482.
% verified: CrossRef 10.1006/aima.1998.1797 2026-10-25
\bibitem{BSY99} M. Balazard, E. Saias and M. Yor. Notes sur la fonction $\zeta$ de Riemann, 2. \emph{Advances in Mathematics} 143 (1999), 284--287.
% verified: CrossRef 10.4213/tvp7 2026-10-24 (Russian original 52 (2007) 111--128 matched; English pages not confirmed, so the page range was removed)
\bibitem{ButuceaTsybakov08a} C. Butucea and A. B. Tsybakov. Sharp optimality in density deconvolution with dominating bias, I. \emph{Theory of Probability and Its Applications} 52 (2008).
% verified: CrossRef 10.1137/s0040585x97982992 2026-10-24
\bibitem{ButuceaTsybakov08b} C. Butucea and A. B. Tsybakov. Sharp optimality in density deconvolution with dominating bias, II. \emph{Theory of Probability and Its Applications} 52 (2008), 237--249.
% verified: CrossRef 10.1073/pnas.41.5.312 2026-09-26
\bibitem{Beurling55} A. Beurling. A closure problem related to the Riemann zeta-function. \emph{Proceedings of the National Academy of Sciences} 41 (1955), 312--314.
% verified: CrossRef 10.1512/iumj.1997.46.1467 2026-10-09 (CrossRef and OpenAlex give no pages; pages 705 to 718 from the printed page numbers of the article PDF)
\bibitem{Biane97} P. Biane. On the free convolution with a semi-circular distribution. \emph{Indiana University Mathematics Journal} 46 (1997), 705--718.
% verified: CrossRef 10.1007/978-3-642-01763-6_7 2026-10-09 (pages 171 to 185; the Lecture Notes volume number is not in CrossRef or OpenAlex and is omitted); arXiv 0811.1490, text checked
\bibitem{Biane09} P. Biane. Matrix valued Brownian motion and a paper by P\'olya. In \emph{S\'eminaire de Probabilit\'es XLII}, Lecture Notes in Mathematics, Springer, 2009, 171--185.
% verified: CrossRef 10.1090/s0273-0979-01-00912-0 2026-09-26
\bibitem{BPY01} P. Biane, J. Pitman and M. Yor. Probability laws related to the Jacobi theta and Riemann zeta functions, and Brownian excursions. \emph{Bulletin of the American Mathematical Society} 38 (2001), 435--465.
% verified: CrossRef 10.1017/cbo9780511721434 2026-10-05 (Theorem 4.12.9 is de Bruijn's exponential Tauberian theorem, as cited by Aurzada and Dereich, Ann. Probab. 37 (2009), arXiv:0805.1330, reference [4]; the 1989 paperback keeps the numbering)
\bibitem{BGT87} N. H. Bingham, C. M. Goldie and J. L. Teugels. \emph{Regular Variation}. Encyclopedia of Mathematics and its Applications 27, Cambridge University Press, 1987.
% verified: CrossRef 10.1016/j.aim.2020.107240, arXiv 2004.11501 2026-11-02
\bibitem{BDV20} F. Broucke, G. Debruyne and J. Vindas. Beurling integers with RH and large oscillation. \emph{Advances in Mathematics} 370 (2020), 107240.
% verified: CrossRef 10.1103/PhysRevLett.77.1472 2026-11-02 (fields only; the content cited was not read)
\bibitem{BogomolnyKeating96} E. B. Bogomolny and J. P. Keating. Gutzwiller's trace formula and spectral statistics: beyond the diagonal approximation. \emph{Physical Review Letters} 77 (1996), 1472--1475.
% verified: CrossRef 10.1088/0305-4470/39/34/010 2026-09-27
\bibitem{BBLM06} E. Bogomolny, O. Bohigas, P. Leboeuf and A. G. Monastra. On the spacing distribution of the Riemann zeros: corrections to the asymptotic result. \emph{Journal of Physics A: Mathematical and General} 39 (2006), 10743--10754.
% verified: CrossRef 10.1006/jnth.1999.2392 2026-09-26
\bibitem{BL99} E. Bombieri and J. C. Lagarias. Complements to Li's criterion for the Riemann hypothesis. \emph{Journal of Number Theory} 77 (1999), 274--287.
% verified: zbMATH 1008.11034 2026-09-26
\bibitem{BombieriStepanov} E. Bombieri. Counting points on curves over finite fields (d'apr\`es S. A. Stepanov). In \emph{S\'eminaire Bourbaki}, expos\'e 430 (1972/73), Lecture Notes in Mathematics 383, Springer, 1974, 234--241. % verified: CrossRef 10.1007/bfb0057311 and zbMATH (Exp. No. 430, LNM 383, 1974) 2026-11-10
\bibitem{Bombieri} E. Bombieri. Remarks on Weil's quadratic functional in the theory of prime numbers, I. \emph{Rendiconti Lincei. Matematica e Applicazioni} (9) 11 (2000), 183--233.
% verified: CrossRef 10.1007/978-1-4612-2948-3 2026-09-26
\bibitem{Bondesson} L. Bondesson. \emph{Generalized Gamma Convolutions and Related Classes of Distributions and Densities}. Lecture Notes in Statistics 76, Springer, 1992.
% verified: CrossRef 10.1090/S0025-5718-1979-0537983-2 2026-09-27
\bibitem{Brent79} R. P. Brent. On the zeros of the Riemann zeta function in the critical strip. \emph{Mathematics of Computation} 33 (1979), 1361--1372.
% verified: CrossRef 10.1007/s00041-013-9292-3 2026-10-24
\bibitem{CandesFG13} E. J. Cand\`es and C. Fernandez-Granda. Super-resolution from noisy data. \emph{Journal of Fourier Analysis and Applications} 19 (2013), 1229--1254.
% verified: OpenAlex 10.1002/cpa.21455 2026-10-24
\bibitem{CandesFG14} E. J. Cand\`es and C. Fernandez-Granda. Towards a mathematical theory of super-resolution. \emph{Communications on Pure and Applied Mathematics} 67 (2014), 906--956.
% verified: CrossRef 10.1090/s0002-9939-04-07677-4 2026-09-25
\bibitem{Cardon2005} D. A. Cardon. Fourier transforms having only real zeros. \emph{Proceedings of the American Mathematical Society} 133 (2005), 1349--1356.
% verified: CrossRef 10.1080/01621459.1988.10478718 2026-10-24
\bibitem{CarrollHall88} R. J. Carroll and P. Hall. Optimal rates of convergence for deconvolving a density. \emph{Journal of the American Statistical Association} 83 (1988), 1184--1186.
% verified: CrossRef 10.1007/s11222-011-9231-6 2026-10-25
\bibitem{CDFG12} F. C\'erou, P. Del Moral, T. Furon and A. Guyader. Sequential Monte Carlo for rare event estimation. \emph{Statistics and Computing} 22 (2012), 795--808.
% verified: CrossRef 10.1214/10-aihp358 2026-10-25
\bibitem{CDG11} F. C\'erou, P. Del Moral and A. Guyader. A nonasymptotic theorem for unnormalized Feynman Kac particle models. \emph{Annales de l'Institut Henri Poincar\'e, Probabilit\'es et Statistiques} 47 (2011), 629--649.
% verified: CrossRef 10.1214/10-aap714 2026-09-27 (pages 1663 to 1693 from zbMATH 1244.60078 and arXiv:1112.5305; CrossRef gives no pages)
\bibitem{CCCS11} X. Chen, L. Cheng, J. Chadam and D. Saunders. Existence and uniqueness of solutions to the inverse boundary crossing problem for diffusions. \emph{Annals of Applied Probability} 21 (2011), 1663--1693.
% verified: CrossRef 10.1214/aos/1176324464 2026-10-24
\bibitem{Chen95} J. Chen. Optimal rate of convergence for finite mixture models. \emph{Annals of Statistics} 23 (1995), 221--233.
% verified: CrossRef 10.1093/biomet/89.3.539 2026-10-25
\bibitem{Chopin02} N. Chopin. A sequential particle filter method for static models. \emph{Biometrika} 89 (2002), 539--551.
% verified: CrossRef 10.1093/biomet/asq021 2026-10-25
\bibitem{ChopinRobert10} N. Chopin and C. P. Robert. Properties of nested sampling. \emph{Biometrika} 97 (2010), 741--755.
% verified: CrossRef 10.1007/bf02385832, 2026-10-28
\bibitem{Chung76} K. L. Chung. Excursions in Brownian motion. \emph{Arkiv f\"or Matematik} 14 (1976), 155--177.
% New keys of editor E27 (ch/c27.tex, cg1 to cg3). Fields are those of the source bibliography, nothing added.
% verified: pending
\bibitem{ComtetMajumdar05} A. Comtet and S. N. Majumdar. Precise asymptotics for a random walker's maximum. \emph{Journal of Statistical Mechanics} (2005), P06013; arXiv cond-mat/0506195.
% verified: CrossRef 10.1007/s00029-021-00689-4 2026-10-17; zbMATH 1480.46084 agrees on volume 27, issue 4, year 2021; article number added as Art. 77 (Paper No. 77 in zbMATH); title comma after formula kept, but CrossRef, zbMATH and Semantic Scholar omit it while arXiv 2006.13771 has it (Springer page not reachable)
\bibitem{ConnesConsani21} A. Connes and C. Consani. Weil positivity and trace formula, the archimedean place. \emph{Selecta Mathematica (N.S.)} 27 (2021), Art. 77, arXiv:2006.13771.
% verified: arXiv 2511.22755 2026-10-17; authors and title match the arXiv API (v1, 2025-11-27); a published version exists, EMS Press chapter CrossRef 10.4171/elm/37/3, pages 39 to 76, 2026-06-23; entry left as the preprint
\bibitem{CCM25} A. Connes, C. Consani and H. Moscovici. Zeta spectral triples. arXiv:2511.22755, 2025.
% verified: zbMATH 0966.11039 2026-09-26
\bibitem{CL00} J. B. Conrey and X.-J. Li. A note on some positivity conditions related to zeta and $L$-functions. \emph{International Mathematics Research Notices} 2000 (2000), 929--940.
% verified: CrossRef 10.2307/3315430 2026-10-24
\bibitem{CP91} B. P. Carlin and N. G. Polson. Inference for nonconjugate Bayesian models using the Gibbs sampler. \emph{Canadian Journal of Statistics} 19 (1991), 399--405.
% verified: CrossRef 10.1007/bf01482072 2026-09-26
\bibitem{Cramer} H. Cram\'er. Ein Mittelwertsatz in der Primzahltheorie. \emph{Mathematische Zeitschrift} 12 (1922), 147--153.
% verified: CrossRef 10.1090/s0002-9947-1986-0846596-4 2026-09-27
\bibitem{CNV86} G. Csordas, T. S. Norfolk and R. S. Varga. The Riemann hypothesis and the Tur\'an inequalities. \emph{Transactions of the American Mathematical Society} 296 (1986), 521--541.
% verified: zbMATH 0807.11059 2026-09-26
\bibitem{COSV93} G. Csordas, A. M. Odlyzko, W. Smith and R. S. Varga. A new Lehmer pair of zeros and a new lower bound for the de Bruijn-Newman constant $\Lambda$. \emph{Electronic Transactions on Numerical Analysis} 1 (1993), 104--111.
% verified: CrossRef 10.1007/bf01205170 2026-09-26
\bibitem{CravenCsordas89} T. Craven and G. Csordas. Jensen polynomials and the Tur\'an and Laguerre inequalities. \emph{Pacific Journal of Mathematics} 136 (1989), 241--260. % verified: DOI 10.2140/pjm.1989.136.241 2026-11-09
\bibitem{CSV94} G. Csordas, W. Smith and R. S. Varga. Lehmer pairs of zeros, the de Bruijn-Newman constant $\Lambda$, and the Riemann hypothesis. \emph{Constructive Approximation} 10 (1994), 107--129.
% verified: pending (not found: CrossRef and OpenAlex searched by title)
\bibitem{CT84} I. Csisz\'ar and G. Tusn\'ady. Information geometry and alternating minimization procedures. \emph{Statistics and Decisions}, Supplement Issue 1 (1984), 205--237.
% verified: CrossRef 10.1007/bf02788653 2026-10-23
\bibitem{CurrySchoenberg66} H. B. Curry and I. J. Schoenberg. On P\'olya frequency functions IV: the fundamental spline functions and their limits. \emph{Journal d'Analyse Math\'ematique} 17 (1966), 71--107.
% verified: zbMATH 1002.11001 2026-09-26
\bibitem{Davenport} H. Davenport. \emph{Multiplicative Number Theory}, third edition, revised by H. L. Montgomery. Graduate Texts in Mathematics 74, Springer, 2000.
% verified: CrossRef 10.1111/j.1467-9868.2006.00553.x 2026-10-25
\bibitem{DDJ06} P. Del Moral, A. Doucet and A. Jasra. Sequential Monte Carlo samplers. \emph{Journal of the Royal Statistical Society, Series B} 68 (2006), 411--436.
% verified: CrossRef 10.1016/j.jmaa.2012.05.011 2026-10-24
\bibitem{deCastroGamboa12} Y. de Castro and F. Gamboa. Exact reconstruction using Beurling minimal extrapolation. \emph{Journal of Mathematical Analysis and Applications} 395 (2012), 336--354.
% verified: pending
\bibitem{deHaanStadtmuller96} L. de Haan and U. Stadtm\"uller. Generalized regular variation of second order. \emph{Journal of the Australian Mathematical Society} 61 (1996), 381--395.
% New bibliography entries of Chapter 34, part b (cf:sec:conc, cf:sec:u34), to be merged into ch/bib.tex by the main session. Fields are those of the source note's reference list; none invented. DelMoral04 is also needed by the other parts of Chapter 34 (add once).
% verified: CrossRef 10.1007/978-1-4684-9393-1 2026-10-25
\bibitem{DelMoral04} P. Del Moral. \emph{Feynman Kac Formulae: Genealogical and Interacting Particle Systems with Applications}. Springer, 2004.
% verified: CrossRef 10.1016/j.spl.2009.07.028 2026-10-23 (Statist. Probab. Lett. 79, 2251-2259; OpenAlex agrees)
\bibitem{DeVoreLorentz93} R. A. DeVore and G. G. Lorentz. \emph{Constructive Approximation}. Grundlehren der mathematischen Wissenschaften 303, Springer, Berlin, 1993. % verified: CrossRef (book DOI, Chapter 9 Approximation by Operators, pp. 267 to 302), Springer ISBN 978-3-540-50627-0, 2026-11-09
\bibitem{Devroye09} L. Devroye. On exact simulation algorithms for some distributions related to Jacobi theta functions. \emph{Statistics and Probability Letters} 79 (2009), 2251--2259.
% verified: CrossRef 10.1112/jlms/s1-11.3.181 (part I, s1-11 no. 3, 181-185, July 1936) and 10.1112/jlms/s1-11.4.307 (part II, s1-11 no. 4, 307-312, October 1936), 2026-10-11; text not seen, so the part that contains the example of period 5 is unchecked and both are cited
\bibitem{DH36} H. Davenport and H. Heilbronn. On the zeros of certain Dirichlet series I, II. \emph{Journal of the London Mathematical Society} 11 (1936), 181--185 and 307--312.
% verified: CrossRef 10.1215/s0012-7094-50-01720-0 2026-09-26; pages 197 to 226 from zbMATH 0038.23302, 2026-10-02
\bibitem{deBruijn} N. G. de Bruijn. The roots of trigonometric integrals. \emph{Duke Mathematical Journal} 17 (1950), 197--226.
% verified: zbMATH 0287.14001 2026-09-26
\bibitem{Deligne74} P. Deligne. La conjecture de Weil. I. \emph{Publications Math\'ematiques de l'IH\'ES} 43 (1974), 273--307.
% verified: CrossRef 10.1007/978-1-4613-8643-8 2026-09-26
\bibitem{Devroye86} L. Devroye. \emph{Non-Uniform Random Variate Generation}. Springer, 1986.
% verified: pending
\bibitem{DL11} D. K. Dimitrov and F. R. Lucas. Higher order Tur\'an inequalities for the Riemann $\xi$ function. \emph{Proceedings of the American Mathematical Society} 139 (2011), 1013--1022.
% verified: CrossRef 10.1090/s0025-5718-1967-0228165-4 2026-10-25
\bibitem{DL67} L. M. Delves and J. N. Lyness. A numerical method for locating the zeros of an analytic function. \emph{Mathematics of Computation} 21 (1967), 543--560.
% verified: CrossRef 10.1007/bfb0103798 2026-10-25
\bibitem{DM00} P. Del Moral and L. Miclo. Branching and interacting particle systems approximations of Feynman Kac formulae with applications to nonlinear filtering. \emph{S\'eminaire de Probabilit\'es XXXIV}, Lecture Notes in Mathematics 1729, 2000, 1--145.
% verified: CrossRef 10.4064/aa200603-23-7 2026-10-08
\bibitem{Dobner21} A. Dobner. A proof of Newman's conjecture for the extended Selberg class. \emph{Acta Arithmetica} 201 (2021), 29--62.
% verified: CrossRef 10.1007/978-3-642-65755-9 2026-10-17; zbMATH 0278.30004 agrees (author, title, Grundlehren series 207, Springer, 1974)
\bibitem{Donoghue74} W. F. Donoghue, Jr. \emph{Monotone Matrix Functions and Analytic Continuation}. Grundlehren der mathematischen Wissenschaften 207, Springer, 1974.
% verified: CrossRef 10.1137/0523074 2026-10-24
\bibitem{Donoho92} D. L. Donoho. Superresolution via sparsity constraints. \emph{SIAM Journal on Mathematical Analysis} 23 (1992), 1309--1331.
% verified: CrossRef 10.1007/s10107-021-01618-1 2026-10-24
\bibitem{DTAB22} R.-A. Dragomir, A. B. Taylor, A. d'Aspremont and J. Bolte. Optimal complexity and certification of Bregman first-order methods. \emph{Mathematical Programming} 194 (2022), 41--83.
% verified: OpenAlex (title match, no DOI) 2026-10-25
\bibitem{Duren70} P. L. Duren. \emph{Theory of $H^p$ Spaces}. Academic Press, 1970.
% verified: pending
\bibitem{Edrei53} A. Edrei. On the generating functions of totally positive sequences II. \emph{Journal d'Analyse Math\'ematique} 2 (1952), 104--109.
% verified: zbMATH 0315.10035 2026-09-26
\bibitem{Edwards} H. M. Edwards. \emph{Riemann's Zeta Function}. Academic Press, 1974.
% verified: CrossRef 10.1002/9780470316658 2026-10-25
\bibitem{EK86} S. N. Ethier and T. G. Kurtz. \emph{Markov Processes: Characterization and Convergence}. Wiley, 1986.
% verified: CrossRef 10.1007/bf01608495 2026-10-01
\bibitem{EMN76} R. S. Ellis, J. L. Monroe and C. M. Newman. The GHS and other correlation inequalities for a class of even ferromagnets. \emph{Communications in Mathematical Physics} 46 (1976), 167--182.
% verified: CrossRef 10.1214/aos/1176348248 2026-10-24
\bibitem{Fan91} J. Fan. On the optimal rates of convergence for nonparametric deconvolution problems. \emph{Annals of Statistics} 19 (1991), 1257--1272.
% verified: CrossRef 10.1103/PhysRevLett.126.056402 2026-10-24
\bibitem{FeiYehGull21} J. Fei, C.-N. Yeh and E. Gull. Nevanlinna analytical continuation. \emph{Physical Review Letters} 126 (2021), 056402.
% verified: OpenAlex 10.1137/18M120885X 2026-10-24
\bibitem{FINUFFT19} A. H. Barnett, J. Magland and L. af Klinteberg. A parallel nonuniform fast Fourier transform library based on an ``exponential of semicircle'' kernel. \emph{SIAM Journal on Scientific Computing} 41 (2019), C479--C504.
% verified: CrossRef 10.2307/2002508 2026-10-24
\bibitem{FL50} G. E. Forsythe and R. A. Leibler. Matrix inversion by a Monte Carlo method. \emph{Mathematical Tables and Other Aids to Computation} 4 (1950), 127--129.
% verified: CrossRef 10.1111/j.1467-9868.2011.01010.x 2026-10-25
\bibitem{FP12} P. Fearnhead and D. Prangle. Constructing summary statistics for approximate Bayesian computation: semi-automatic approximate Bayesian computation. \emph{Journal of the Royal Statistical Society, Series B} 74 (2012), 419--474.
% verified: CrossRef 10.53846/goediss-5113 2026-09-25
\bibitem{Gabcke79} W. Gabcke. \emph{Neue Herleitung und explizite Restabsch\"atzung der Riemann-Siegel-Formel}. Dissertation, Georg-August-Universit\"at G\"ottingen, 1979.
% verified: CrossRef 10.1214/16-EJS1119 2026-10-24
\bibitem{GaoVdV16} F. Gao and A. W. van der Vaart. Posterior contraction rates for deconvolution of Dirichlet Laplace mixtures. \emph{Electronic Journal of Statistics} 10 (2016), 608--627.
% verified: zbMATH 0136.11201 2026-10-17; title, authors, and Academic Press, New York and London, 1964 (XIV, 384 p.) agree; OpenAlex 1964 record agrees
\bibitem{GelfandVilenkin4} I. M. Gel'fand and N. Ya. Vilenkin. \emph{Generalized Functions, Volume 4: Applications of Harmonic Analysis}. Academic Press, 1964.
% verified: CrossRef 10.1214/ss/1177011137 2026-10-24
\bibitem{Geyer92} C. J. Geyer. Practical Markov chain Monte Carlo. \emph{Statistical Science} 7 (1992), 473--483.
% New bibliography entries for the merge of note N (reconcile-N.md section 5), checked by web-researcher on 26 October 2026; merged into ch/bib.tex by the main session.
% verified: CrossRef 10.1215/s0012-7094-07-13634-2 2026-10-26 (journal title; the arXiv title math/0511182 adds "formula")
\bibitem{GHK07} S. M. Gonek, C. P. Hughes and J. P. Keating. A hybrid Euler-Hadamard product for the Riemann zeta function. \emph{Duke Mathematical Journal} 136 (2007), no. 3, 507--549. doi:10.1215/s0012-7094-07-13634-2, arXiv:math/0511182.
% verified: CrossRef 10.1017/s0305004100042225 2026-10-09
\bibitem{Goldie67} C. Goldie. A class of infinitely divisible random variables. \emph{Proceedings of the Cambridge Philosophical Society} 63 (1967), 1141--1143.
% verified: pending
\bibitem{GoldstonMontgomery87} D. A. Goldston and H. L. Montgomery. Pair correlation of zeros and primes in short intervals. In \emph{Analytic Number Theory and Diophantine Problems}, Birkh\"auser, 1987, 183--203.
% verified: zbMATH 1297.11101 2026-09-26
\bibitem{GonekMontgomery13} S. M. Gonek and H. L. Montgomery. Zeros of a family of approximations of the Riemann zeta-function. \emph{International Mathematics Research Notices} 2013 (2013), 4712--4733.
% verified: pending (fields as the source note gives them, JASA 75, 42; the end page and the title were not checked)
\bibitem{GoodGaskins80} I. J. Good and R. A. Gaskins. Density estimation and bump hunting by the penalized likelihood method. \emph{Journal of the American Statistical Association} 75 (1980), 42.
% verified: publisher URL http://numbers.computation.free.fr/Constants/Miscellaneous/zetazeros1e13-1e24.pdf 2026-09-27
\bibitem{Gourdon04} X. Gourdon. The $10^{13}$ first zeros of the Riemann Zeta function, and zeros computation at very large height. Manuscript, 2004, \url{http://numbers.computation.free.fr/Constants/Miscellaneous/zetazeros1e13-1e24.pdf}, accessed 27 September 2026.
% verified: CrossRef 10.1007/bf02421310 2026-09-27
\bibitem{Gram03} J.-P. Gram. Note sur les z\'eros de la fonction $\zeta(s)$ de Riemann. \emph{Acta Mathematica} 27 (1903), 289--304.
% verified: CrossRef 10.1073/pnas.1902572116 2026-10-01
\bibitem{GORZ19} M. Griffin, K. Ono, L. Rolen and D. Zagier. Jensen polynomials for the Riemann zeta function and other sequences. \emph{Proceedings of the National Academy of Sciences} 116 (2019), 11103--11110.
% verified: CrossRef 10.1007/bf02465533 2026-10-12 (English translation of Liet. Mat. Rink. 39 (1999), 40 to 51; the publisher's page returned a client challenge, so the paper itself was not read)
\bibitem{Grigelionis99} B. Grigelionis. Processes of Meixner type. \emph{Lithuanian Mathematical Journal} 39 (1999), 33--41.
% verified: pending
\bibitem{Guinand48} A. P. Guinand. A summation formula in the theory of prime numbers. \emph{Proceedings of the London Mathematical Society} 50 (1948), 107--119.
% verified: CrossRef 10.4007/annals.2026.203.2.6 2026-09-27 (pages 623 to 675 from the journal; CrossRef gives no pages)
\bibitem{GM26} L. Guth and J. Maynard. New large value estimates for Dirichlet polynomials. \emph{Annals of Mathematics} 203 (2026), 623--675.
% verified: pending (not found: CrossRef and OpenAlex searched by title; book)
\bibitem{Hadamard23} J. Hadamard. \emph{Lectures on Cauchy's Problem in Linear Partial Differential Equations}. Yale University Press, 1923.
% verified: arXiv 0910.5228 2026-10-16 (title and authors checked on the abstract page)
\bibitem{Haglund09} J. Haglund. Some conjectures on the zeros of approximates to the Riemann $\Xi$-function and incomplete gamma functions. arXiv:0910.5228, 2009.
% verified: zbMATH 45.0716.04 2026-10-22 (no DOI; title, author, C. R. 158 (1914), 1012-1014 confirmed)
\bibitem{Hardy14} G. H. Hardy. Sur les z\'eros de la fonction $\zeta(s)$ de Riemann. \emph{Comptes Rendus de l'Acad\'emie des Sciences Paris} 158 (1914), 1012--1014.
% verified: CrossRef 10.1007/bf02789198 2026-10-22 (J. Anal. Math. 55 (1990), 59-95 confirmed)
\bibitem{Hejhal90} D. A. Hejhal. On a result of G. P\'olya concerning the Riemann $\xi$-function. \emph{Journal d'Analyse Math\'ematique} 55 (1990), 59--95.
% verified: OpenAlex 10.4007/annals.2011.174.2.4 2026-10-24
\bibitem{Hiary11} G. A. Hiary. Fast methods to compute the Riemann zeta function. \emph{Annals of Mathematics} 174 (2011).
% verified: zbMATH 1336.11061 2026-10-22 (no DOI found; Indag. Math., new series 27, no. 2 (2016), 524-533 confirmed)
\bibitem{Hiary16} G. A. Hiary. An explicit van der Corput estimate for $\zeta(1/2+it)$. \emph{Indagationes Mathematicae} 27 (2016), 524--533.
% verified: zbMATH 1159.11001 2026-09-26
\bibitem{boHardyWright} G. H. Hardy and E. M. Wright. \emph{An Introduction to the Theory of Numbers}, sixth edition, revised by D. R. Heath-Brown and J. H. Silverman. Oxford University Press, 2008.
% verified: CrossRef 10.1016/j.jnt.2021.06.032 2026-10-11 (title, journal, volume 235, pages 219 to 241, issued June 2022, authors Hasanalizade, Shen, Wong); arXiv 2107.06506v1 read for (1.7), (2.7), Corollaries 1.2 and 1.4 and footnote 2 (the error in Trudgian 2014)
\bibitem{HSW22} E. Hasanalizade, Q. Shen and P.-J. Wong. Counting zeros of the Riemann zeta function. \emph{Journal of Number Theory} 235 (2022), 219--241.
% verified: CrossRef 10.1007/bf01565437 2026-10-09
\bibitem{Hecke36} E. Hecke. \"Uber die Bestimmung Dirichletscher Reihen durch ihre Funktionalgleichung. \emph{Mathematische Annalen} 112 (1936), 664--699.
% verified: OpenAlex 10.1109/29.56027 2026-10-24
\bibitem{HuaSarkar90} Y. Hua and T. K. Sarkar. Matrix pencil method for estimating parameters of exponentially damped undamped sinusoids in noise. \emph{IEEE Transactions on Acoustics, Speech, and Signal Processing} 38 (1990), 814--824.
% verified: CrossRef 10.1016/0370-1573(95)00074-7 2026-10-24
\bibitem{JarrellGubernatis96} M. Jarrell and J. E. Gubernatis. Bayesian inference and the analytic continuation of imaginary-time quantum Monte Carlo data. \emph{Physics Reports} 269 (1996), 133--195.
% New bibliography entries of ch/cf3.tex (Chapter 34, prefix cf:), to be merged into ch/bib.tex by the main session. Existing keys used: Selberg46, HSW22, PT21, DH36, Spira94, Vinogradov58. Fields of Kadiri05 and MT15 are the source note's list; the others are from the reconcile-L.md list and memory, every one to be checked.
% verified: CrossRef 10.4064/aa117-4-1 2026-10-25
\bibitem{Kadiri05} H. Kadiri. Une r\'egion explicite sans z\'eros pour la fonction $\zeta$ de Riemann. \emph{Acta Arithmetica} 117 (2005), 303--339.
% verified: OpenAlex 10.1007/978-1-4757-4015-8 2026-10-24
\bibitem{Kallenberg02} O. Kallenberg. \emph{Foundations of Modern Probability}, 2nd ed. Springer, 2002.
% E3 (ht5), keys used in ch/ht5.tex; entries below not yet in ch/bib.tex
% verified: pending
\bibitem{Karlin68} S. Karlin. \emph{Total Positivity}, Volume I. Stanford University Press, 1968.
% verified: OpenAlex title match, no DOI in record (NeurIPS 28 (2015) 2845--2853) 2026-10-24
\bibitem{KBB15} W. Krichene, A. Bayen and P. L. Bartlett. Accelerated mirror descent in continuous and discrete time. \emph{Advances in Neural Information Processing Systems} 28 (2015).
% verified: CrossRef 10.1017/s0021900200094468, 2026-10-28
\bibitem{Kennedy76} D. P. Kennedy. The distribution of the maximum Brownian excursion. \emph{Journal of Applied Probability} 13 (1976), 371--376.
% verified: CrossRef 10.1214/11-aos920 2026-10-24
\bibitem{KnapikVZ11} B. T. Knapik, A. W. van der Vaart and J. H. van Zanten. Bayesian inverse problems with Gaussian priors. \emph{Annals of Statistics} 39 (2011), 2626--2657.
% verified: CrossRef 10.1080/03610926.2012.681417 2026-10-24
\bibitem{KnapikVZ13} B. T. Knapik, A. W. van der Vaart and J. H. van Zanten. Bayesian recovery of the initial condition for the heat equation. \emph{Communications in Statistics, Theory and Methods} 42 (2013), 1294--1313.
\bibitem{Lagarias06} J. C. Lagarias. Hilbert spaces of entire functions and Dirichlet $L$-functions. In \emph{Frontiers in Number Theory, Physics, and Geometry I}, Springer, 2006, 365--377.
% verified: CrossRef 10.1007/s00020-013-2091-z 2026-11-02 (metadata only; paper not read)
\bibitem{KreinLanger14} M. G. Krein and H. Langer. Continuation of Hermitian positive definite functions and related questions. \emph{Integral Equations and Operator Theory} 78 (2014), 1--69.
% verified: arXiv 2405.04869 2026-10-24 (author initials, title and year corrected from arXiv)
\bibitem{Leong24} N. Leong. Explicit estimates for the logarithmic derivative and the reciprocal of the Riemann zeta function. Preprint arXiv:2405.04869, 2024.
% verified: OpenAlex 10.1090/mmono/150 2026-10-23 (book record)
\bibitem{Levin96} B. Ya. Levin. \emph{Lectures on Entire Functions}. American Mathematical Society, 1996.
% verified: CrossRef 10.1016/j.acha.2020.10.004 2026-10-24
\bibitem{LiLiao21} W. Li and W. Liao. Stable super resolution limit and smallest singular value of restricted Fourier matrices. \emph{Applied and Computational Harmonic Analysis} 51 (2021), 118--156.
% verified: CrossRef 10.1093/biomet/81.1.27 2026-10-24
\bibitem{LiuWongKong94} J. S. Liu, W. H. Wong and A. Kong. Covariance structure of the Gibbs sampler with applications to the comparisons of estimators and augmentation schemes. \emph{Biometrika} 81 (1994), 27--40.
% verified: CrossRef 10.1086/111605 2026-10-24
\bibitem{Lucy74} L. B. Lucy. An iterative technique for the rectification of observed distributions. \emph{Astronomical Journal} 79 (1974), 745--754.
% verified: CrossRef 10.1002/rsa.20135 2026-10-24
\bibitem{LV07} L. Lov\'asz and S. Vempala. The geometry of logconcave functions and sampling algorithms. \emph{Random Structures and Algorithms} 30 (2007), 307--358.
% verified: CrossRef 10.1007/978-3-540-87557-4 2026-10-24
\bibitem{Meister09} A. Meister. \emph{Deconvolution Problems in Nonparametric Statistics}. Springer, 2009.
% verified: zbMATH 0044.07402 2026-10-22 (title, author, year 1951 and pages 61-70 confirmed; zbMATH gives the series as "Ser. A 2" where the chapter has "Ser. A 1 (1951), no. 2"; volume not settled, editor to check)
\bibitem{Miyatake51a} O. Miyatake. Note on Riemann's $\xi$ function. \emph{J. Inst. Polytech. Osaka City Univ. Ser. A} 1 (1951), no. 2, 61--70.
% verified: zbMATH 0044.07401 2026-10-22 (title, author, year 1951 and pages 39-59 confirmed; zbMATH gives the series as "Ser. A 2" where the chapter has "Ser. A 1 (1951), no. 2"; volume not settled, editor to check)
\bibitem{Miyatake51b} O. Miyatake. On the distribution of zero points of a function which is related to Riemann's $\xi$ function. \emph{J. Inst. Polytech. Osaka City Univ. Ser. A} 1 (1951), no. 2, 39--59.
% verified: pending
\bibitem{MontgomerySoundararajan04} H. L. Montgomery and K. Soundararajan. Primes in short intervals. \emph{Communications in Mathematical Physics} 252 (2004), 589--617.
% verified: CrossRef 10.4064/aa166-2-4 2026-10-25
\bibitem{MSZ14} Yu. Matiyasevich, F. Saidak and P. Zvengrowski. Horizontal monotonicity of the modulus of the zeta function, Dirichlet $L$-functions, and related functions. \emph{Acta Arithmetica} 166 (2014), 189--200.
% verified: CrossRef 10.1016/j.jnt.2015.05.010 2026-10-25
\bibitem{MT15} M. J. Mossinghoff and T. S. Trudgian. Nonnegative trigonometric polynomials and a zero free region for the Riemann zeta function. \emph{Journal of Number Theory} 157 (2015), 329--349.
% verified: CrossRef 10.1007/s40993-023-00498-y and arXiv 2212.06867 2026-10-25 (journal, volume and year from the record)
\bibitem{MTY} M. J. Mossinghoff, T. S. Trudgian and A. Yang. Explicit zero free regions for the Riemann zeta function. \emph{Research in Number Theory} 10 (2024). doi:10.1007/s40993-023-00498-y. arXiv:2212.06867.
\bibitem{Nesterov83} Y. Nesterov. A method of solving a convex programming problem with convergence rate $O(1/k^2)$. \emph{Soviet Mathematics Doklady} 27 (1983), 372--376.
% verified: OpenAlex 10.1214/12-AOS1065 2026-10-24
\bibitem{Nguyen13} X. Nguyen. Convergence of latent mixing measures in finite and infinite mixture models. \emph{Annals of Statistics} 41 (2013), 370--400.
% verified: CrossRef 10.1007/s10208-013-9150-3 2026-10-24
\bibitem{OC15} B. O'Donoghue and E. J. Cand\`es. Adaptive restart for accelerated gradient schemes. \emph{Foundations of Computational Mathematics} 15 (2015), 715--732.
% verified: CrossRef 10.1090/s0002-9947-1988-0961614-2 2026-10-24
\bibitem{OdlyzkoSchonhage88} A. M. Odlyzko and A. Sch\"onhage. Fast algorithms for multiple evaluations of the Riemann zeta function. \emph{Transactions of the American Mathematical Society} 309 (1988), 797--809.
% verified: manuscript (title page of the note, misc/reviews/nick-bspin/source)
\bibitem{P26bspin} N. G. Polson. Bessel spin ferromagnets: Lee Yang level repulsion, the GHS inequality, and a moment criterion for the Riemann Hypothesis. Manuscript, October 2026.
% New bibliography entries of Chapter 35, part a (ch/c35.tex, ch/bsy1.tex), to be merged into ch/bib.tex by the main session. Existing keys used: Titchmarsh, Levin96, PT21, HSW22, Davenport, DH36, Spira94, BSO07, OAI26a. BSY99 is defined in bib-new-cfa.tex (Chapter 34 part a) and is used here unchanged. P26bsy: title and authors from the source; Duren70: from reconcile-M.md section 3 (not in the source); nothing else added.
% verified: manuscript, author's note 2026-10-25
\bibitem{P26bsy} N. G. Polson and V. Sokolov. Balazard, Saias and Yor meet Thorin: Cauchy averages and their Monte Carlo solution. Manuscript, October 2026.
% New bibliography entries of Chapter 34, part a (ch/c34.tex, ch/cf1.tex), to be merged into ch/bib.tex by the main session. Existing keys used: PSH626, PS23, Lagarias, SD10, P17v5 (not used in the end), DH36, Titchmarsh, Spira94, PT21, Vinogradov58, Kallenberg02, OAI26a. BSY99 is shared with Chapter 35 (add once). Fields from the editor's knowledge of the papers, none invented; all pending.
% verified: manuscript, author's note 2026-10-25
\bibitem{P26cf} N. G. Polson. Cauchy deconvolution from $U_1$ to $U_{1/2}$: particle filtering and generative Bayesian computation along the Cauchy flow of the Riemann xi function. Manuscript, October 2026.
% verified: pending
\bibitem{P26cg} N. G. Polson. The Comtet and Majumdar identity, the $2\sin^2$ lemma, and approximating families for the reciprocal Xi function. Manuscript, October 2026.
% verified: pending (the manuscript is the source of this chapter, untracked, see misc/reviews/nick-lvj_sech/source5)
\bibitem{P26es} N. G. Polson. The explicit formula as a deconvolution. Manuscript, October 2026.
% verified: pending (title and sole author from the source file; defined here only if the editor of ch/c25.tex has not added it)
\bibitem{P26evt} N. G. Polson. Extreme values, hitting times and a path to RH. Manuscript, October 2026.
% verified: manuscript (title page of the note, misc/reviews/nick-fchap/source3, dated October 28, 2026)
\bibitem{P26excursion} N. G. Polson. The excursion maximum behind the pole cancellation: work on the open problems. Manuscript, October 2026.
% verified: manuscript (chapter title of the note, misc/reviews/nick-fchap/source; the note has no title page or date)
\bibitem{P26fchap} N. G. Polson. Ferromagnets and the Riemann zeros. Manuscript, October 2026.
% verified: manuscript (title page of the note, misc/reviews/nick-gumbel/source2)
\bibitem{P26gumbel} N. G. Polson. A Gumbel mixture path to the Riemann Hypothesis. Manuscript, October 2026.
% verified: manuscript (title page of the note, misc/reviews/nick-gumbel/source2)
\bibitem{P26gumbelop} N. G. Polson. The Gamma skeleton, heat flow, and the arithmetic displacement: work on the open problems of the Gumbel mixture path. Manuscript, October 2026.
% verified: pending
\bibitem{P26heat} N. G. Polson. The natural EGGC family for the reciprocal Xi function: heat flow, the $2\sin^2$ lemma, $z$ distributions and a class on $(0,\infty)$. Manuscript, October 2026.
% verified: pending (source note of the second half of Chapter 28; the note is headed "Progress on the Open Problems", a note for the book's editor; no journal)
\bibitem{P26hermite} N. G. Polson. Progress on the open problems. Manuscript, October 2026.
% verified: pending
\bibitem{P26ht} N. G. Polson. Hitting times and the reciprocal Xi function. Manuscript, October 2026.
% verified: manuscript (title page of the note lattice_bias_law.tex, dated October 28, 2026)
\bibitem{P26lbias} N. G. Polson. How close can a finite approximant get? The lattice bias law for Dyck zeros and the universal Lehmer obstruction. Manuscript, October 2026.
% verified: pending
\bibitem{P26lo} N. G. Polson. The Thorin measure at level one and the backward Cauchy flow. Manuscript, October 2026.
% verified: manuscript (title page of the note, misc/reviews/nick-gumbel/source2)
\bibitem{P26path} N. G. Polson. A path to the Riemann Hypothesis: ferromagnetic skeleton, Coulomb gas, and the one missing step. Manuscript, October 2026.
% verified: pending
\bibitem{P26pg3} N. G. Polson. The Xi P\'olya Gamma law III: self decomposability without RH, Li's coefficients as cumulants, the motion of zeros, and Fej\'er mixtures. Manuscript, October 2026.
% verified: manuscript (title page of the note, misc/reviews/nick-fchap/source2, dated October 28, 2026)
\bibitem{P26pole} N. G. Polson. Pole cancellation: why no component of Riemann's mixture is ferromagnetic, and what a ferromagnetic solution must therefore be. Manuscript, October 2026.
% verified: pending (source note of the first half of Chapter 28; title and date from the note, no journal)
\bibitem{P26trunc} N. G. Polson. Thorin truncation and the zeros of $L(s,\chi_4)$. Manuscript, October 2026.
% verified: manuscript (title page of the note two_uniform_spins.tex, dated October 28, 2026)
\bibitem{P26tus} N. G. Polson. Two uniform spins: the exact ferromagnet behind $\xi$, Riemann's contour proof as a Lee Yang trace, and the transform RH asks about. Manuscript, October 2026.
% verified: manuscript (title page of the note, misc/reviews/nick-wferro/source)
\bibitem{P26wferro} N. G. Polson. Williams, ferromagnetic random variables and the Riemann Xi function: a ferromagnetic family and its reciprocal. Manuscript, October 2026.
% verified: pending
\bibitem{P26wt} N. G. Polson. Weil positivity as the Thorin condition. Manuscript, October 2026.
% verified: CrossRef 10.1214/18-EJS1430 2026-11-08
\bibitem{PadillaPolsonScott18} O. H. Madrid Padilla, N. G. Polson and J. G. Scott. A deconvolution path for mixtures. \emph{Electronic Journal of Statistics} 12 (2018), 1717--1751.
% verified: CrossRef 10.1214/aos/1176343003 2026-10-23 (Ann. Statist. 3, 119-131; OpenAlex agrees)
\bibitem{Pickands75} J. Pickands. Statistical inference using extreme order statistics. \emph{Annals of Statistics} 3 (1975), 119--131.
% verified: pending
\bibitem{PolsonEuler26} N. G. Polson. The Euler product of $1/\zeta$ as a L\'evy measure, the $2\sin^2$ lemma, and the reciprocal Xi function as a generalized gamma convolution. Note, October 2026.
% verified: pending
\bibitem{PolsonResidue26} N. G. Polson. A residue expansion of the reciprocal Xi function and generalized gamma convolution methods. Note, October 2026.
% verified: OpenAlex arXiv 1803.09138 2026-10-24 (arXiv record; NeurIPS venue not separately confirmed)
\bibitem{PolsonRockova18} N. G. Polson and V. Ro\v{c}kov\'a. Posterior concentration for sparse deep learning. \emph{Advances in Neural Information Processing Systems} 31 (2018).
% verified: the note itself, title from its first page (misc/reviews/nick-rxe/source, N:41) 2026-10-26
\bibitem{PolsonRXi26} N. G. Polson. The reciprocal Xi function as the Laplace transform of an EGGC: mixing laws, the Euler product, Selberg and Barnes. Note, October 2026.
% verified: crossref 10.1111/j.1467-9868.2011.01015.x 2026-10-27
\bibitem{PolsonScott12} N. G. Polson and J. G. Scott. Local shrinkage rules, L\'evy processes and regularized regression. \emph{Journal of the Royal Statistical Society, Series B} 74 (2012), 287--311.
% verified: CrossRef 10.1111/rssb.12130 2026-11-08
\bibitem{PolsonScott16} N. G. Polson and J. G. Scott. Mixtures, envelopes and hierarchical duality. \emph{Journal of the Royal Statistical Society, Series B} 78 (2016), 701--727.
% verified: CrossRef 10.3390/e27070683 and arXiv 2305.14972 2026-10-24
\bibitem{PS23} N. G. Polson and V. Sokolov. Generative AI for Bayesian computation. \emph{Entropy} 27 (2025), 683. arXiv:2305.14972.
% verified: pending (not found: CrossRef and OpenAlex searched by title; only SSRN 10.2139/ssrn.7160360, a tutorial with a different title)
\bibitem{PSGBC26} N. G. Polson and V. Sokolov. \emph{Generative Bayesian Computation: Quantile Neural Networks for Inference and Surrogates}. Chapman and Hall/CRC, 2026.
% verified: pending (not found: only SSRN 10.2139/ssrn.7192018 under a longer title)
\bibitem{PSH626} N. G. Polson and V. Sokolov. \emph{Hilbert's Sixth Problem}. KDP, 2026.
% verified: CrossRef 10.1080/01621459.2013.829001 2026-10-22 (J. Amer. Statist. Assoc. 108 (2013), 1339-1349 confirmed; authors Polson, Scott, Windle)
\bibitem{PSW13} N. G. Polson, J. G. Scott and J. Windle. Bayesian Inference for Logistic Models Using P\'olya--Gamma Latent Variables. \emph{Journal of the American Statistical Association} 108 (2013), 1339--1349.
% verified: pending (not found: software citation, not indexed by CrossRef or OpenAlex)
\bibitem{pythonflint} The python-flint developers. \emph{python-flint: Python bindings for FLINT and Arb}. Software, \url{https://github.com/flintlib/python-flint}.
% verified: pending
\bibitem{Rademacher59} H. Rademacher. On the Phragm\'en--Lindel\"of theorem and some applications. \emph{Mathematische Zeitschrift} 72 (1959), 192--204.
% verified: CrossRef 10.1007/bf02127580 2026-10-23
\bibitem{Renyi53} A. R\'enyi. On the theory of order statistics. \emph{Acta Mathematica Academiae Scientiarum Hungaricae} 4 (1953), 191--231.
% verified: CrossRef 10.1214/aoap/1034625254 2026-10-23
\bibitem{RGG97} G. O. Roberts, A. Gelman and W. R. Gilks. Weak convergence and optimal scaling of random walk Metropolis algorithms. \emph{Annals of Applied Probability} 7 (1997), 110--120.
% verified: CrossRef 10.1364/JOSA.62.000055 2026-10-24
\bibitem{Richardson72} W. H. Richardson. Bayesian based iterative method of image restoration. \emph{Journal of the Optical Society of America} 62 (1972), 55--59.
% verified: CrossRef 10.2307/2005479 2026-10-25
\bibitem{RS75} J. B. Rosser and L. Schoenfeld. Sharper bounds for the Chebyshev functions $\theta(x)$ and $\psi(x)$. \emph{Mathematics of Computation} 29 (1975), 243--269.
% verified: CrossRef 10.1017/s096249291000005x 2026-10-24
\bibitem{Rump10} S. M. Rump. Verification methods: rigorous results using floating point arithmetic. \emph{Acta Numerica} 19 (2010), 287--449.
% verified: pending (not found: JMLR 17 (2016) not indexed; OpenAlex has only NeurIPS 2014 and arXiv 1503.01243)
\bibitem{SBC16} W. Su, S. Boyd and E. J. Cand\`es. A differential equation for modeling Nesterov's accelerated gradient method: theory and insights. \emph{Journal of Machine Learning Research} 17 (2016), 1--43.
% verified: CrossRef 10.1093/imaiai/iax006 2026-10-24
\bibitem{SchiebingerRR18} G. Schiebinger, E. Robeva and B. Recht. Superresolution without separation. \emph{Information and Inference} 7 (2018), 1--30.
% verified: CrossRef 10.1109/tmi.1982.4307558 2026-10-24
\bibitem{ShVa82} L. A. Shepp and Y. Vardi. Maximum likelihood reconstruction for emission tomography. \emph{IEEE Transactions on Medical Imaging} 1 (1982), 113--122.
% New bibliography entries of Chapter 34, part d (ch/cf4.tex), to be merged into ch/bib.tex by the main session. Fields are those of the source's reference list (L:1722 to 1752); none invented (no page ranges). Existing keys used: BSO07, BBLM06, DH36, PT21, Spira94. The source's key is in the comment.
% verified: CrossRef 10.1214/06-BA127 2026-10-25
\bibitem{Skilling06} J. Skilling. Nested sampling for general Bayesian computation. \emph{Bayesian Analysis} 1 (2006), 833--859.
% verified: pending
\bibitem{Spitzer56} F. Spitzer. A combinatorial lemma and its application to probability theory. \emph{Transactions of the American Mathematical Society} 82 (1956), 323--339.
% verified: issue listing geodesic.mathdoc.fr/item/MZM_1970_8_4 (Mathnet mzm9606), Mat. Zametki 8:4 (1970) 419-429, 2026-10-25
\bibitem{Stechkin70} S. B. Stechkin. Zeros of the Riemann zeta function. \emph{Matematicheskie Zametki} 8 (1970), 419--429.
% verified: CrossRef 10.1080/02331889008802238 2026-10-24 (title corrected to Deconvolving)
\bibitem{StefanskiCarroll90} L. A. Stefanski and R. J. Carroll. Deconvolving kernel density estimators. \emph{Statistics} 21 (1990), 169--184.
% verified: CrossRef 10.1090/s0002-9947-1925-1501297-5 2026-09-27
\bibitem{Hutchinson25} J. I. Hutchinson. On the roots of the Riemann zeta function. \emph{Transactions of the American Mathematical Society} 27 (1925), 49--60.
% verified: zbMATH 0006.39701 2026-10-17 (author, title, Cambridge Tracts 30, Cambridge University Press, 1932; no CrossRef record found; zbMATH 0715.11045 gives the same series number 30 for the 1990 Cambridge Mathematical Library reprint)
\bibitem{Ingham32} A. E. Ingham. \emph{The Distribution of Prime Numbers}. Cambridge Tracts in Mathematics and Mathematical Physics 30, Cambridge University Press, 1932.
% verified: zbMATH 0025.02704 2026-09-26
\bibitem{Ingham40} A. E. Ingham. On the estimation of $N(\sigma,T)$. \emph{Quarterly Journal of Mathematics} 11 (1940), 291--292.
% verified: zbMATH 1059.11001 2026-09-27
\bibitem{IwaniecKowalski04} H. Iwaniec and E. Kowalski. \emph{Analytic Number Theory}. American Mathematical Society Colloquium Publications 53, American Mathematical Society, 2004.
% verified: CrossRef 10.1214/07-ps118 2026-09-26; pages 346 to 415 from zbMATH 1189.60035, 2026-10-02
\bibitem{JRY} L. F. James, B. Roynette and M. Yor. Generalized gamma convolutions, Dirichlet means, Thorin measures, with explicit examples. \emph{Probability Surveys} 5 (2008), 346--415.
% verified: CrossRef 10.1109/tc.2017.2690633 2026-09-26
\bibitem{Arb} F. Johansson. Arb: efficient arbitrary-precision midpoint-radius interval arithmetic. \emph{IEEE Transactions on Computers} 66 (2017), 1281--1292.
% verified: CrossRef 10.1215/kjm/1250522571 2026-10-17; zbMATH 0421.40009 and OpenAlex agree; 1978 journal name kept (now Kyoto Journal of Mathematics)
\bibitem{Kasahara78} Y. Kasahara. Tauberian theorems of exponential type. \emph{Journal of Mathematics of Kyoto University} 18 (1978), 209--219.
\bibitem{Keiper92} J. B. Keiper. Power series expansions of Riemann's $\xi$ function. \emph{Mathematics of Computation} 58 (1992), 765--773.
% verified: CrossRef 10.1007/bf00538895 2026-09-26
\bibitem{Kent80} J. T. Kent. Eigenvalue expansions for diffusion hitting times. \emph{Zeitschrift f\"ur Wahrscheinlichkeitstheorie und Verwandte Gebiete} 52 (1980), 309--319.
% verified: CrossRef 10.1214/aop/1176993924 2026-10-09 (CrossRef gives no pages; pages 207 to 219 from zbMATH 0483.60075; statement checked in the Project Euclid abstract)
\bibitem{Kent82} J. T. Kent. The spectral decomposition of a diffusion hitting time. \emph{Annals of Probability} 10 (1982), 207--219.
% verified: CrossRef 10.1016/j.aim.2009.04.003 2026-09-29
\bibitem{KKL09} H. Ki, Y.-O. Kim and J. Lee. On the de Bruijn--Newman constant. \emph{Advances in Mathematics} 222 (2009), 281--306.
% verified: CrossRef 10.5802/aif.3600 2026-09-27
\bibitem{KPS24} T. Konstantopoulos, P. Patie and R. Sarkar. A new class of solutions to the van Dantzig problem, the Lee--Yang property, and the Riemann hypothesis. \emph{Annales de l'Institut Fourier} 74 (2024), 377--421.
% verified: zbMATH 0086.03803 2026-10-17 (author, journal, volume 13, no. 4(82), pages 185 to 192, 1958; the zbMATH title is German, Abschaetzungen trigonometrischer Summen und ihre Anwendungen; no CrossRef or OpenAlex record found; mathnet.ru unreachable)
\bibitem{Korobov58} N. M. Korobov. Estimates of trigonometric sums and their applications. \emph{Uspekhi Matematicheskikh Nauk} 13 (1958), no. 4, 185--192.
% verified: CrossRef 10.4064/aa-89-3-217-234 2026-09-26
\bibitem{Lagarias} J. C. Lagarias. On a positivity property of the Riemann $\xi$-function. \emph{Acta Arithmetica} 89 (1999), 217--234.
% verified: CrossRef 10.4064/aa120-2-4 2026-10-01 (statement checked in the arXiv abstract, math/0601653)
\bibitem{Lagarias05} J. C. Lagarias. Zero spacing distributions for differenced $L$-functions. \emph{Acta Arithmetica} 120 (2005), 159--184.
% verified: CrossRef 10.1007/bf01456808 2026-10-17; zbMATH 43.0332.01 agrees on author, title, journal, volume and pages; year corrected from 1911 to 1912 (CrossRef and zbMATH both date it 1912)
\bibitem{Landau11} E. Landau. \"Uber die Nullstellen der Zetafunktion. \emph{Mathematische Annalen} 71 (1912), 548--564.
% verified: CrossRef 10.1002/j.1538-7305.1962.tb03279.x 2026-10-17; zbMATH 0184.08603 agrees on author, title, journal, volume, pages and year (title differs only in capitalization and hyphenation)
\bibitem{LandauPollak62} H. J. Landau and H. O. Pollak. Prolate spheroidal wave functions, Fourier analysis and uncertainty, III: the dimension of the space of essentially time and band limited signals. \emph{Bell System Technical Journal} 41 (1962), 1295--1336.
% verified: CrossRef 10.1016/0022-247x(80)90241-3 2026-10-17; zbMATH 0471.47029 agrees on author, title, journal, volume, pages and year
\bibitem{LandauWidom80} H. J. Landau and H. Widom. Eigenvalue distribution of time and frequency limiting. \emph{Journal of Mathematical Analysis and Applications} 77 (1980), 469--481.
% verified: zbMATH 0860.65028 2026-09-26
\bibitem{LawsonHanson74} C. L. Lawson and R. J. Hanson. \emph{Solving Least Squares Problems}. Prentice Hall, 1974.
% verified: CrossRef 10.1103/physrev.87.410 2026-09-26
\bibitem{LeeYang} T. D. Lee and C. N. Yang. Statistical theory of equations of state and phase transitions. II. Lattice gas and Ising model. \emph{Physical Review} 87 (1952), 410--419.
% verified: CrossRef 10.1112/plms/s3-20.2.303 2026-09-26
\bibitem{Lehman70} R. S. Lehman. On the distribution of zeros of the Riemann zeta-function. \emph{Proceedings of the London Mathematical Society} (3) 20 (1970), 303--320.
% verified: CrossRef 10.1007/bf02401102 2026-09-27
\bibitem{Lehmer56} D. H. Lehmer. On the roots of the Riemann zeta-function. \emph{Acta Mathematica} 95 (1956), 291--298.
% verified: CrossRef 10.1112/s0025579300001753 2026-09-27
\bibitem{Lehmer56b} D. H. Lehmer. Extended computation of the Riemann zeta-function. \emph{Mathematika} 3 (1956), 102--108.
% verified: zbMATH 0152.06703 2026-10-16 (book; series and year checked)
\bibitem{Levin64} B. Ja. Levin. \emph{Distribution of Zeros of Entire Functions}. Translations of Mathematical Monographs 5, American Mathematical Society, 1964.
% verified: CrossRef 10.1006/jnth.1997.2137 2026-09-26
\bibitem{Li97} X.-J. Li. The positivity of a sequence of numbers and the Riemann hypothesis. \emph{Journal of Number Theory} 65 (1997), 325--333.
% verified: CrossRef 10.1007/bf01213009 2026-10-01
\bibitem{LiebSokal81} E. H. Lieb and A. D. Sokal. A general Lee-Yang theorem for one-component and multicomponent ferromagnets. \emph{Communications in Mathematical Physics} 80 (1981), 153--179.
% verified: CrossRef 10.1007/bf01170633 2026-10-17; zbMATH 0008.11301 agrees (author, title, Math. Z. 38, pages 177 to 216, 1934)
\bibitem{Lowner34} K. L\"owner. \"Uber monotone Matrixfunktionen. \emph{Mathematische Zeitschrift} 38 (1934), 177--216.
\bibitem{Liu26a} B. Liu. Improved zero-free half-planes for the quasi-Riemann hypothesis. arXiv:2610.12234 [math.NT], submitted 8 October 2026, v2 12 October 2026. Preprint, not refereed. % verified: arXiv listing, reference check of Nick Polson of 21 October 2026
\bibitem{Liu26b} B. Liu. Zero-free half-planes for Hecke $L$-functions over number fields. arXiv:2610.26835, submitted 19 October 2026. Preprint, not refereed. % verified: arXiv listing, reference check of Nick Polson of 21 October 2026
% verified: CrossRef 10.1111/j.2517-6161.1993.tb01925.x 2026-10-24 (pages corrected to 569--598)
\bibitem{VardiLee93} Y. Vardi and D. Lee. From image deblurring to optimal investments: maximum likelihood solutions for positive linear inverse problems. \emph{Journal of the Royal Statistical Society, Series B} 55 (1993), 569--598.
% verified: CrossRef 10.1214/aop/1176995155, 2026-10-28
\bibitem{Vervaat79} W. Vervaat. A relation between Brownian bridge and Brownian excursion. \emph{The Annals of Probability} 7 (1979), 143--149.
% New bibliography entries of Chapter 35, part b (ch/bsy2.tex), to be merged into ch/bib.tex by the main session. Existing keys used: Titchmarsh, Spira94, BSO07, DH36, PT21, HSW22. BSY99 is in bib-new-cfa.tex (shared). Duren70 and any Levin96 page are in part a.
% verified: arXiv 1409.3601 2026-10-25
\bibitem{VLMC14} N. G. Polson and J. G. Scott. Vertical likelihood Monte Carlo. arXiv:1409.3601.
% verified: pending
\bibitem{vonKoch01} H. von Koch. Sur la distribution des nombres premiers. \emph{Acta Mathematica} 24 (1901), 159--182.
% verified: pending
\bibitem{Voros10} A. Voros. \emph{Zeta Functions over Zeros of Zeta Functions}. Lecture Notes of the Unione Matematica Italiana 8, Springer, 2010.
% verified: CrossRef 10.1111/j.1467-9469.2007.00585.x 2026-10-24
\bibitem{VR08} R. Varadhan and C. Roland. Simple and globally convergent methods for accelerating the convergence of any EM algorithm. \emph{Scandinavian Journal of Statistics} 35 (2008), 335--353.
% verified: CrossRef 10.1080/01621459.1985.10477119 2026-10-24 (pages corrected to 8--20)
\bibitem{VSK85} Y. Vardi, L. A. Shepp and L. Kaufman. A statistical model for positron emission tomography. \emph{Journal of the American Statistical Association} 80 (1985), 8--20.
% verified: arXiv 2610.10584 abstract page 2026-10-25 (sole author Henning Wunderlich, v1 of 6 October 2026)
\bibitem{Wun26} H. Wunderlich. Finite-radius Jensen and relative-entropy formulations of the Balazard-Saias-Yor criterion. arXiv:2610.10584, 2026.
\bibitem{XuSang26} Y. Xu and H. Sang. Amplified moments and a refinement of the quasi-Riemann hypothesis. arXiv:2610.15291, October 2026. Preprint, not refereed. % verified: arXiv listing, reference check of Nick Polson of 21 October 2026
\bibitem{Xylouris11} T. Xylouris. On the least prime in an arithmetic progression and estimates for the zeros of Dirichlet $L$-functions. \emph{Acta Arithmetica} 150 (2011), no.~1, 65--91. doi:10.4064/aa150-1-4. % verified: IMPAN page, reference check of Nick Polson of 21 October 2026; the Linnik constant of this paper is $L=5.2$, the value $5$ is from the later Bonn dissertation
% verified: CrossRef 10.1007/s00440-021-01105-w 2026-10-09; arXiv 1504.00350, text checked (the form (p +_d q)(x) = p^(D) q^(D) x^d after Definition 1.1, the Hermite polynomials H_d = e^{-D^2/2} x^d in Section 1.2, and the remark there that limits of the finite convolutions give the free ones, citing Marcus 2015 and Arizmendi and Perales 2018)
\bibitem{MSS22} A. W. Marcus, D. A. Spielman and N. Srivastava. Finite free convolutions of polynomials. \emph{Probability Theory and Related Fields} 182 (2022), 807--848.
% verified: CrossRef 10.1090/pspum/024/9944 2026-09-27
\bibitem{Montgomery73} H. L. Montgomery. The pair correlation of zeros of the zeta function. In \emph{Analytic Number Theory}, H. G. Diamond, ed., Proceedings of Symposia in Pure Mathematics 24, American Mathematical Society, 1973, 181--193.
% verified: CrossRef 10.1007/978-3-0348-5438-2_42 2026-10-14; the record's book title omits the subtitle 'To the Memory of Paul Turan'
\bibitem{Montgomery83} H. L. Montgomery. Zeros of approximations to the zeta function. In \emph{Studies in Pure Mathematics: To the Memory of Paul Tur\'an}, Birkh\"auser, Basel, 1983, 497--506.
% verified: PyPI mpmath-1.3.0 2026-09-26
\bibitem{mpmath} The mpmath development team. \emph{mpmath: a Python library for arbitrary-precision floating-point arithmetic}, version 1.3.0, 2023.
% verified: CrossRef 10.1002/cpa.3160270203 2026-09-27
\bibitem{Newman74} C. M. Newman. Zeros of the partition function for generalized Ising systems. \emph{Communications on Pure and Applied Mathematics} 27 (1974), 143--159.
% verified: CrossRef 10.1090/s0002-9939-1976-0434982-5 2026-09-25
\bibitem{Newman1976} C. M. Newman. Fourier transforms with only real zeros. \emph{Proceedings of the American Mathematical Society} 61 (1976), 245--251.
% verified: CrossRef 10.1007/bf01888165 2026-10-01
\bibitem{Newman91} C. M. Newman. The GHS inequality and the Riemann hypothesis. \emph{Constructive Approximation} 7 (1991), 389--399.
% verified: Zenodo concept DOI 10.5281/zenodo.23461586 resolves (checked 2026-10-20 by the authors' assistant): title 'Kernel-Concavity Obstructions and Fourier-Side Laguerre Bounds for the Riemann Xi Programme', author Tony Newton, version v1 of 19 October 2026, a newer version exists
\bibitem{Newton26} T. Newton. Kernel-concavity obstructions and Fourier-side Laguerre bounds for the Riemann $\Xi$ programme. Zenodo, 2026, \url{https://doi.org/10.5281/zenodo.23461586}, accessed 20 October 2026.
% verified: CrossRef 10.1214/ecp.v14-1488 2026-09-26 (pages 396 to 411 from the journal; CrossRef gives no pages)
\bibitem{NY09} A. Nikeghbali and M. Yor. The Barnes $G$ function and its relations with sums and products of generalized gamma convolution variables. \emph{Electronic Communications in Probability} 14 (2009), 396--411.
% verified: zbMATH 0037.35401 2026-09-26
\bibitem{Nyman50} B. Nyman. \emph{On the One-Dimensional Translation Group and Semi-Group in Certain Function Spaces}. Thesis, University of Uppsala, 1950.
% verified: CrossRef 10.1090/s0025-5718-1987-0866115-0 2026-09-25
\bibitem{Odlyzko87} A. M. Odlyzko. On the distribution of spacings between zeros of the zeta function. \emph{Mathematics of Computation} 48 (1987), 273--308.
% verified: CrossRef 10.1090/s0002-9947-1988-0961614-2 2026-09-26
\bibitem{OS88} A. M. Odlyzko and A. Sch\"onhage. Fast algorithms for multiple evaluations of the Riemann zeta function. \emph{Transactions of the American Mathematical Society} 309 (1988), 797--809.
% verified: URL https://www-users.cse.umn.edu/~odlyzko/unpublished/zeta.10to20.1992.pdf 2026-09-27
\bibitem{Odlyzko92} A. M. Odlyzko. \emph{The $10^{20}$-th zero of the Riemann zeta function and 175 million of its neighbors}. Unpublished manuscript, 1992, \url{https://www-users.cse.umn.edu/~odlyzko/unpublished/zeta.10to20.1992.pdf}, accessed 27 September 2026.
% verified: CrossRef 10.1023/a:1016677511798 2026-09-29 (the bound -2.63e-9 is computed after (3.14) in Section 3; the abstract and (1.12) state -2.7e-9)
\bibitem{Odlyzko00} A. M. Odlyzko. An improved bound for the de Bruijn--Newman constant. \emph{Numerical Algorithms} 25 (2000), 293--303.
% verified: CrossRef 10.1090/conm/290/04578 2026-09-27
\bibitem{Odlyzko01} A. M. Odlyzko. The $10^{22}$-nd zero of the Riemann zeta function. In \emph{Dynamical, Spectral, and Arithmetic Zeta Functions}, M. L. Lapidus and M. van Frankenhuysen, eds., Contemporary Mathematics 290, American Mathematical Society, 2001, 139--144.
% verified: URL https://www-users.cse.umn.edu/~odlyzko/zeta_tables/ 2026-09-27
\bibitem{OdlyzkoTables} A. M. Odlyzko. \emph{Tables of zeros of the Riemann zeta function}. Online, \url{https://www-users.cse.umn.edu/~odlyzko/zeta_tables/}, accessed 27 September 2026.
% verified: publisher URL https://dlmf.nist.gov/help/cite 2026-09-27
\bibitem{DLMF} F. W. J. Olver, A. B. Olde Daalhuis, D. W. Lozier, B. I. Schneider, R. F. Boisvert, C. W. Clark, B. R. Miller, B. V. Saunders, H. S. Cohl and M. A. McClain, eds. \emph{NIST Digital Library of Mathematical Functions}. \url{https://dlmf.nist.gov/}, Release 1.2.8 of 2026-09-15.
\bibitem{OAI26a} OpenAI. The quasi-Riemann hypothesis: a zero-free half-plane $\re(s)>7/8$. Preprint dated 30 September 2026, released 6 October 2026 with a formalisation in Lean, \url{https://github.com/openai/math} (family 003), accessed 7 October 2026.
\bibitem{OAI26b} OpenAI. The quasi-Riemann hypothesis. Preprint dated 5 October 2026, released 6 October 2026, \url{https://github.com/openai/math} (family 003), accessed 7 October 2026.
% verified: CrossRef 10.1216/jiea/1181074930 2026-09-26; pages 397 to 423 from zbMATH 1044.60076, 2026-10-02
\bibitem{Peskir02} G. Peskir. On integral equations arising in the first-passage problem for Brownian motion. \emph{Journal of Integral Equations and Applications} 14 (2002), 397--423.
% verified: CrossRef 10.4153/cjm-2003-014-x 2026-10-12; PDF read at Cambridge Core (formulas (1), (4), (67) and Table 5 as cited)
\bibitem{PitmanYor03} J. Pitman and M. Yor. Infinitely divisible laws associated with hyperbolic functions. \emph{Canadian Journal of Mathematics} 55 (2003), 292--330.
% verified: zbMATH 1345.11064 2026-09-26
\bibitem{Platt16} D. J. Platt. Numerical computations concerning the GRH. \emph{Mathematics of Computation} 85 (2016), 3009--3027.
% verified: CrossRef 10.1112/blms.12460 2026-09-26
\bibitem{PT21} D. J. Platt and T. S. Trudgian. The Riemann hypothesis is true up to $3\cdot10^{12}$. \emph{Bulletin of the London Mathematical Society} 53 (2021), 792--797.
% verified: arXiv 1708.02653v24 2026-09-27
\bibitem{P17} N. G. Polson. On Hilbert's 8th problem. arXiv:1708.02653v24, 2026.
% verified: arXiv 1708.02653v5 2026-09-27
\bibitem{P17v5} N. G. Polson. On Hilbert's 8th problem. arXiv:1708.02653v5, 2017.
% verified: publisher URL https://www.amazon.com/dp/B0HKPL1X1S 2026-09-27 (Independently published, 22 September 2026, 488 pages, ISBN 9798176411898; title and subtitle from the title page of the interior PDF, whose copyright page reads "First edition, 2026" and prints no publisher or ISBN; pinpoints follow this edition, not the draft of about 195 pages)
\bibitem{P26} N. G. Polson. \emph{On Hilbert's 8th Problem: the reciprocal Xi clock, Thorin measures, van Dantzig pairs, zero free approximations, hitting times and the reciprocal zeta function}, first edition. Independently published, 2026.
% verified: manuscript N/A 2026-09-27 (title, sole author and date from the source file eta_paper.tex; no arXiv record)
\bibitem{P26eta} N. G. Polson. Thorin measures of the Dedekind eta function. Manuscript, September 2026.
% verified: manuscript N/A 2026-10-14 (title from the source file chapter_max_wald_thorin.tex, received on 14 October 2026 as chapter_max_wald.zip; sole author as for the first author's other manuscripts; no arXiv record)
\bibitem{P26mw} N. G. Polson. The P\'olya density as a maximum order statistic and the Thorin measure of the reciprocal Xi function. Manuscript, October 2026.
% verified: manuscript N/A 2026-09-27 (title, byline and date from reciprocal-zeta/paper/reciprocal_zeta_thorin.tex)
\bibitem{P26rz} N. G. Polson. Thorin type representations of the reciprocal zeta function on the critical line. Manuscript, September 2026, now Part IX of the first edition of \emph{On Hilbert's 8th Problem} (2026).
% verified: manuscript N/A 2026-10-09 (title, sole author and date from the source files wigner_free_thorin_v3.tex and wigner_free_thorin_v7.tex, the latter dated October 2026 (version 7); no arXiv record)
\bibitem{P26wig} N. G. Polson. Wigner's semicircle, free probability and the Thorin framework for the reciprocal Xi function. Manuscript, version 7, October 2026.
% verified: manuscript (title page of the note, misc/reviews/nick-pgVIII/source)
\bibitem{P26pg8} N. G. Polson. The Xi P\'olya Gamma law VIII: why the exponent is $\cos(2\pi/9)$: pole cancellation, a Meixner convolution, and the structure of the P\'olya mixture. Manuscript, October 2026.
% verified: manuscript N/A 2026-10-27 (title page of the note, misc/reviews/nick-pgX/source; no arXiv record)
\bibitem{P26pg10} N. G. Polson. The Xi P\'olya Gamma law X: interlacing families: the cosh family $\operatorname{Re}\xi(\sigma+it)$, the variance shrink theorem, and the zero free region as the extent of an interlacing family. Manuscript, October 2026.
% verified: manuscript N/A 2026-09-27 (authors' source thorin_xi_computation.tex, title page dated September 2026; not on arXiv)
\bibitem{PS26} N. G. Polson and V. Sokolov. Computing the Thorin measure of the reciprocal Xi function: a P\'olya route and a total variation comparison with the Riemann atoms. Manuscript, September 2026.
% verified: manuscript N/A 2026-10-14 (working report of both authors, dated 7 October 2026, cited by the source files of Chapter 14 received on 14 October 2026 as crh_ch14.zip, which supersede it; no arXiv record)
\bibitem{FA26} N. G. Polson and V. Sokolov. Four attacks on the computational path: a working report. Manuscript, 7 October 2026.
% verified: manuscript N/A 2026-10-14 (draft book of both authors, cited for its Chapter 21 by the source files of Chapter 14 received on 14 October 2026 as crh_ch14.zip; no published record)
\bibitem{LG26} N. G. Polson and V. Sokolov. \emph{On Langlands Conjectures: Probabilistic Methods}, Chapter 21. Draft, 2026.
% verified: CrossRef 10.1007/bf02565336 2026-09-26
\bibitem{Polya} G. P\'olya. Bemerkung \"uber die Integraldarstellung der Riemannschen $\xi$-Funktion. \emph{Acta Mathematica} 48 (1926), 305--317.
% verified: CrossRef 10.1515/crll.1927.158.6 2026-10-16
\bibitem{Polya27} G. P\'olya. \"Uber trigonometrische Integrale mit nur reellen Nullstellen. \emph{Journal f\"ur die reine und angewandte Mathematik} 158 (1927), 6--18.
% verified: CrossRef 10.1007/s40687-019-0193-1 2026-09-26
\bibitem{Polymath} D. H. J. Polymath. Effective approximation of heat flow evolution of the Riemann $\xi$ function, and a new upper bound for the de Bruijn--Newman constant. \emph{Research in the Mathematical Sciences} 6 (2019), paper 31.
% verified: zbMATH 0895.30001 2026-09-27
\bibitem{Remmert98} R. Remmert. \emph{Classical Topics in Complex Function Theory}, translated by L. Kay. Graduate Texts in Mathematics 172, Springer, 1998.
% verified: CrossRef 10.1214/aoap/1177004773 2026-09-26; pages 446 to 460 from zbMATH 0836.60087, 2026-10-02
\bibitem{RS95} G. O. Roberts and C. F. Shortland. The hazard rate tangent approximation for boundary hitting times. \emph{Annals of Applied Probability} 5 (1995), 446--460.
% verified: CrossRef 10.1017/fmp.2020.6 2026-09-26
\bibitem{RodgersTao} B. Rodgers and T. Tao. The de Bruijn--Newman constant is non-negative. \emph{Forum of Mathematics, Pi} 8 (2020), e6.
% verified: CrossRef 10.1215/ijm/1255631807 2026-09-28 (pages 64 to 94 from zbMATH 0122.05001; CrossRef gives no pages)
\bibitem{RS62} J. B. Rosser and L. Schoenfeld. Approximate formulas for some functions of prime numbers. \emph{Illinois Journal of Mathematics} 6 (1962), 64--94.
% verified: zbMATH 0191.44504 2026-09-27 (editor and publisher from Open Library and from Brent, Math. Comp. 33 (1979), reference 23); pages 70 to 76 from zbMATH 0191.44504, 2026-10-02
\bibitem{RYS69} J. B. Rosser, J. M. Yohe and L. Schoenfeld. Rigorous computation and the zeros of the Riemann zeta-function. In \emph{Information Processing 68: Proceedings of IFIP Congress 1968}, volume 1, A. J. H. Morrell, ed., North-Holland, 1969, 70--76.
% verified: zbMATH 0792.11034 2026-09-26; pages 415 to 440 from zbMATH 0792.11034, which adds the supplement S17 to S23, 2026-10-02
\bibitem{Rumely93} R. Rumely. Numerical computations concerning the ERH. \emph{Mathematics of Computation} 61 (1993), 415--440.
% verified: zbMATH 0973.60001 2026-09-26
\bibitem{Sato99} K. Sato. \emph{L\'evy Processes and Infinitely Divisible Distributions}. Cambridge Studies in Advanced Mathematics 68, Cambridge University Press, 1999.
% verified: CrossRef 10.1007/bf02790092 2026-09-27 (accent in Polya from zbMATH 0045.37602 and the original first page; CrossRef omits it)
\bibitem{Schoenberg51} I. J. Schoenberg. On P\'olya frequency functions. I. The totally positive functions and their Laplace transforms. \emph{Journal d'Analyse Math\'ematique} 1 (1951), 331--374.
% verified: CrossRef 10.1090/ulect/033 2026-10-17; zbMATH 1057.30036 agrees on author, title, series, volume 33, publisher and year (title differs only in capitalization)
\bibitem{Seip04} K. Seip. \emph{Interpolation and Sampling in Spaces of Analytic Functions}. University Lecture Series 33, American Mathematical Society, 2004.
% verified: zbMATH 0061.08402 2026-09-26
\bibitem{Selberg46} A. Selberg. Contributions to the theory of the Riemann zeta-function. \emph{Archiv for Mathematik og Naturvidenskab} 48 (1946), 89--155.
% verified: zbMATH 0004.10501 2026-09-26; pages 45 to 80 from zbMATH 0004.10501, 2026-10-02
\bibitem{Siegel32} C. L. Siegel. \"Uber Riemanns Nachla{\ss} zur analytischen Zahlentheorie. \emph{Quellen und Studien zur Geschichte der Mathematik, Astronomie und Physik, Abteilung B: Studien} 2 (1932), 45--80.
% verified: CrossRef 10.1002/j.1538-7305.1961.tb03976.x 2026-10-17; volume 40, issue 1, pages 43 to 63, 1961; zbMATH 0184.08601 and OpenAlex agree
\bibitem{SlepianPollak61} D. Slepian and H. O. Pollak. Prolate spheroidal wave functions, Fourier analysis and uncertainty, I. \emph{Bell System Technical Journal} 40 (1961), 43--63.
% verified: CrossRef 10.1007/s10998-010-1037-3 2026-09-26
\bibitem{SD10} J. Sondow and C. Dumitrescu. A monotonicity property of Riemann's xi function and a reformulation of the Riemann hypothesis. \emph{Periodica Mathematica Hungarica} 60 (2010), 37--40.
% verified: zbMATH 0226.10040 2026-09-26
\bibitem{Spira67} R. Spira. Zeros of approximate functional approximations. \emph{Mathematics of Computation} 21 (1967), 41--48.
% verified: CrossRef 10.1090/s0025-5718-1994-1254148-8 2026-10-11 (Math. Comp. 63, no. 208, first page 747; CrossRef gives no last page); 747--748 as in Ivic, arXiv:math/0311162, ref. [58]
\bibitem{Spira94} R. Spira. Some zeros of the Titchmarsh counterexample. \emph{Mathematics of Computation} 63 (1994), 747--748.
% verified: CrossRef 10.1214/aoms/1177698806 2026-10-09
\bibitem{Steutel67} F. W. Steutel. Note on the infinite divisibility of exponential mixtures. \emph{Annals of Mathematical Statistics} 38 (1967), 1303--1305.
% verified: CrossRef 10.1098/rspa.1936.0192 2026-09-27
\bibitem{Titchmarsh36} E. C. Titchmarsh. The zeros of the Riemann zeta-function. \emph{Proceedings of the Royal Society of London, Series A} 157 (1936), 261--263.
% verified: zbMATH 0601.10026 2026-09-26
\bibitem{Titchmarsh} E. C. Titchmarsh. \emph{The Theory of the Riemann Zeta-Function}, second edition, revised by D. R. Heath-Brown. Oxford University Press, 1986.
% verified: CrossRef 10.1137/130932132 2026-09-26
\bibitem{TW14} L. N. Trefethen and J. A. C. Weideman. The exponentially convergent trapezoidal rule. \emph{SIAM Review} 56 (2014), 385--458.
% verified: CrossRef 10.1090/s0025-5718-2011-02470-1 2026-09-26
\bibitem{Trudgian11} T. S. Trudgian. Improvements to Turing's method. \emph{Mathematics of Computation} 80 (2011), 2259--2279.
% verified: CrossRef 10.4064/aa148-3-2 2026-09-27
\bibitem{Trudgian11b} T. S. Trudgian. On the success and failure of Gram's Law and the Rosser Rule. \emph{Acta Arithmetica} 148 (2011), 225--256.
% verified: CrossRef 10.1016/j.jnt.2013.07.017 2026-09-26
\bibitem{Trudgian} T. S. Trudgian. An improved upper bound for the argument of the Riemann zeta-function on the critical line II. \emph{Journal of Number Theory} 134 (2014), 280--292.
% verified: 2 November 2026 at arxiv.org/abs/2206.03682 (abstract read; the sign convention of the screw function is not in the abstract)
\bibitem{Suzuki22} M. Suzuki. Aspects of the screw function corresponding to the Riemann zeta function. arXiv:2206.03682.
% verified: CrossRef 10.1007/bf01459119 2026-10-11; text not checked (cited through BS18, Theorem 10 (ii), reference [29])
\bibitem{Chebotarev28} N. Tschebotareff. \"Uber die Realit\"at von Nullstellen ganzer transzendenter Funktionen. \emph{Mathematische Annalen} 99 (1928), 660--686.
% verified: CrossRef 10.1112/plms/s3-3.1.99 2026-09-26
\bibitem{Turing53} A. M. Turing. Some calculations of the Riemann zeta-function. \emph{Proceedings of the London Mathematical Society} (3) 3 (1953), 99--117.
% verified: CrossRef 10.1090/s0025-5718-1986-0829637-3 2026-09-27
\bibitem{LRW86} J. van de Lune, H. J. J. te Riele and D. T. Winter. On the zeros of the Riemann zeta function in the critical strip. IV. \emph{Mathematics of Computation} 46 (1986), 667--681.
% verified: zbMATH 0097.26302 2026-10-17 (author, journal, volume 22, pages 161 to 164, 1958; the zbMATH title is German, Eine neue Abschaetzung der Funktion zeta(1+it); no CrossRef record found; mathnet.ru unreachable)
\bibitem{Vinogradov58} I. M. Vinogradov. A new estimate of the function $\zeta(1+it)$. \emph{Izvestiya Akademii Nauk SSSR, Seriya Matematicheskaya} 22 (1958), 161--164.
% verified: zbMATH 0049.03205 2026-09-26; pages 252 to 265 from zbMATH 0049.03205, 2026-10-02
\bibitem{Weil52} A. Weil. Sur les ``formules explicites'' de la th\'eorie des nombres premiers. \emph{Communications du S\'eminaire Math\'ematique de l'Universit\'e de Lund}, tome suppl\'ementaire (1952), 252--265.
% verified: zbMATH 0063.08245 2026-09-27
\bibitem{Widder41} D. V. Widder. \emph{The Laplace Transform}. Princeton Mathematical Series 6, Princeton University Press, 1941.
% verified: zbMATH 0744.60100 2026-09-26
\bibitem{Williams90} D. Williams. Brownian motion and the Riemann zeta-function. In \emph{Disorder in Physical Systems}, G. R. Grimmett and D. J. A. Welsh, eds., Oxford University Press, 1990, 361--372.
% verified: CrossRef 10.1214/aop/1176995474 2026-09-26; pages 523 to 531 from zbMATH 0394.60017, 2026-10-02
\bibitem{Yamazato78} M. Yamazato. Unimodality of infinitely divisible distribution functions of class $L$. \emph{Annals of Probability} 6 (1978), 523--531.
% verified: CrossRef 10.1137/0715088 2026-10-25
\bibitem{YKS78} S. Yakowitz, J. E. Krimmel and F. Szidarovszky. Weighted Monte Carlo integration. \emph{SIAM Journal on Numerical Analysis} 15 (1978), 1289--1300.
% verified: zbMATH 0817.11041 2026-09-27 (DOI 10.2969/aspm/02110281; editors from zbMATH 0771.00036; conference at Tokyo Institute of Technology, August 14 to 18, 1990, per Project Euclid)
\bibitem{Yoshida92} H. Yoshida. On Hermitian forms attached to zeta functions. In \emph{Zeta Functions in Geometry (Tokyo, 1990)}, N. Kurokawa and T. Sunada, eds., Advanced Studies in Pure Mathematics 21, Kinokuniya, 1992, 281--325.
% verified: CrossRef 10.1007/s00208-006-0051-5 2026-11-02
\bibitem{Zhang07} W.-B. Zhang. Beurling primes with RH and Beurling primes with large oscillation. \emph{Mathematische Annalen} 337 (2007), 671--704.
% verified: arXiv 2608.24827 2026-10-17; author Zhu, Xuefeng and title match v2 (2026-09-02) via the arXiv API and zbMATH; v1 (2026-08-25) lists Chuk, Marcus alone with a different title, so the author of record changed between versions (the arXiv comment of v2 says: v1 was a 9 page announcement, author name and affiliation updated); the book cites v2, the full version of 34 pages
\bibitem{Zhu26} X. Zhu. Weil positivity in compact windows: a finite reduction, certified two-sided bounds, and a Landau--Widom decay law. arXiv:2608.24827v2, 2 September 2026.
\bibitem{NCPolymath7} T. Tao. Polymath15, seventh thread: going below 0.48. Blog post, \emph{What's new}, 28 March 2018, \url{https://terrytao.wordpress.com/2018/03/28/polymath15-seventh-thread-going-below-0-48/}. % verified: text read through the TagTeam mirror of the blog feed (dated 2018-03-29 there); author name and the exact wordpress URL inferred from the blog, not shown on the mirror, so the URL is UNCHECKED
\bibitem{NCOAI26c} OpenAI. Uniform exclusion of Landau--Siegel zeros. Preprint dated 1 October 2026, \url{https://github.com/openai/math} (family 003). % verified: CONTENTS.md and lean/docs/003.md of the repository, read 28 October 2026
\bibitem{NCQuanta26} Quanta Magazine. AI proved the quasi-Riemann hypothesis. What happens now? 19 October 2026, \url{https://www.quantamagazine.org/ai-proved-the-quasi-riemann-hypothesis-what-happens-now-20261019}. % verified: article fetched 28 October 2026
\bibitem{NCDunnRadziwill24} A. Dunn and M. Radziwi\l\l. Bias in cubic Gauss sums: Patterson's conjecture. \emph{Annals of Mathematics} (2) 200 (2024), 967--1057. % verified: journal page annals.math.princeton.edu/2024/200-3/p03, DOI 10.4007/annals.2024.200.3.3
\bibitem{NCPatterson77} S. J. Patterson. A cubic analogue of the theta series. \emph{Journal f\"ur die reine und angewandte Mathematik} 296 (1977), 125--161. % verified: reference [38] of OAI26b with DOI 10.1515/crll.1977.296.125; journal page not opened
\bibitem{NCKubota69} T. Kubota. \emph{On Automorphic Functions and the Reciprocity Law in a Number Field}. Lectures in Mathematics, Department of Mathematics, Kyoto University, No.~2, Kinokuniya, Tokyo, 1969. % verified: reference [26] of OAI26b only; UNCHECKED against a library record
\bibitem{NCGoldmakherLouvel13} L. Goldmakher and B. Louvel. A quadratic large sieve inequality over number fields. \emph{Mathematical Proceedings of the Cambridge Philosophical Society} 154 (2013), 193--212. % verified: reference [13] of OAI26b with DOI 10.1017/S0305004112000370; journal page not opened
\bibitem{NCIwaniecConv06} H. Iwaniec. Conversations on the exceptional character. In \emph{Analytic Number Theory} (Cetraro, 2002), Lecture Notes in Mathematics 1891, Springer, 2006, 97--132. % verified: DOI 10.1007/978-3-540-36364-4_3, as supplied by the authors; the text was not read, and we do not know whether it describes the elliptic curve family of Section~1
\bibitem{NCConreyIwaniec02} J. B. Conrey and H. Iwaniec. Spacing of zeros of Hecke $L$-functions and the class number problem. \emph{Acta Arithmetica} 103 (2002), 259--312. % verified: DOI 10.4064/aa103-3-5, as supplied by the authors
\bibitem{NCConrey03} J. B. Conrey. The Riemann hypothesis. \emph{Notices of the American Mathematical Society} 50 (2003), 341--353. % UNCHECKED: fields from Conrey's publication list only; the survey text was not read
\bibitem{NCMestre86} J.-F. Mestre. Formules explicites et minorations de conducteurs de vari\'et\'es alg\'ebriques. \emph{Compositio Mathematica} 58 (1986), 209--232. % verified: Numdam CM_1986__58_2_209_0 and EuDML 89769, 2026-11-04 (no DOI)
\bibitem{NCBrumer92} A. Brumer. The average rank of elliptic curves I. \emph{Inventiones Mathematicae} 109 (1992), 445--472. % verified: CrossRef 10.1007/bf01232033 2026-11-04
\bibitem{NCBoas54} R. P. Boas. \emph{Entire Functions}. Academic Press, New York, 1954. % verified: publisher and year, WorldCat 847696, 2026-11-04; pages 124 to 132 for Akhiezer's factorisation per secondary sources, book not read
\bibitem{NCGrossZagier86} B. H. Gross and D. B. Zagier. Heegner points and derivatives of $L$-series. \emph{Inventiones Mathematicae} 84 (1986), 225--320. % verified: CrossRef 10.1007/BF01388809 2026-11-04
\bibitem{NCKolyvagin89} V. A. Kolyvagin. Finiteness of $E(\mathbb Q)$ and $\mathrm{Sha}(E,\mathbb Q)$ for a subclass of Weil curves. \emph{Mathematics of the USSR Izvestiya} 32 (1989), 523--541. % verified: CrossRef 10.1070/im1989v032n03abeh000779 2026-11-04 (English translation of the 1988 Izvestiya paper)
\bibitem{BohrCourant} H. Bohr and R. Courant. Neue Anwendungen der Theorie der Diophantischen Approximationen auf die Riemannsche Zetafunktion. \emph{Journal f\"ur die reine und angewandte Mathematik} 144 (1914), 249--274. % UNCHECKED: fields from the bibliography of the note, not seen in the original
\bibitem{NCKeiper96} J. B. Keiper. On the zeros of the Ramanujan $\tau$-Dirichlet series in the critical strip. \emph{Mathematics of Computation} 65 (1996), 1613--1619. % UNCHECKED: fields from the author's bibliography; the paper was not read, and the count of 12069 zeros below height 6397 is the author's
\bibitem{NCBerndtDixit10} B. C. Berndt and A. Dixit. A transformation formula involving the Gamma and Riemann zeta functions in Ramanujan's lost notebook. In \emph{The Legacy of Alladi Ramakrishnan in the Mathematical Sciences}, Springer, 2010. % UNCHECKED: fields from the author's bibliography, no pages given; the proof of the formula was not read
\bibitem{NCRamanujanLost88} S. Ramanujan. \emph{The Lost Notebook and Other Unpublished Papers}. Narosa, New Delhi, 1988. % UNCHECKED: fields from the author's bibliography
\bibitem{BookerThorne14} A. R. Booker and F. Thorne. Zeros of $L$-functions outside the critical strip. \emph{Algebra \& Number Theory} 8 (2014), 2027--2042. % verified: CrossRef 10.2140/ant.2014.8.2027 2026-11-10
\bibitem{NCHamburger21} H. Hamburger. \"Uber die Riemannsche Funktionalgleichung der $\zeta$-Funktion. \emph{Mathematische Zeitschrift} 10 (1921), 240--254. % verified: 10.1007/bf01211612 2026-11-09
\bibitem{NCBeurling37} A. Beurling. Analyse de la loi asymptotique de la distribution des nombres premiers g\'en\'eralis\'es. I. \emph{Acta Mathematica} 68 (1937), 255--291. % verified: 10.1007/bf02546666 2026-11-09
\bibitem{NCDMV06} H. G. Diamond, H. L. Montgomery and U. M. A. Vorhauer. Beurling primes with large oscillation. \emph{Mathematische Annalen} 334 (2006), 1--36. % verified: 10.1007/s00208-005-0638-2 2026-11-09
\bibitem{NCDiamondZhang16} H. G. Diamond and W.-B. Zhang. \emph{Beurling Generalized Numbers}. Mathematical Surveys and Monographs 213, American Mathematical Society, 2016. % verified: 10.1090/surv/213 2026-11-09
\bibitem{NCCohnElkies03} H. Cohn and N. Elkies. New upper bounds on sphere packings I. \emph{Annals of Mathematics} 157 (2003), 689--714. % verified: 10.4007/annals.2003.157.689 2026-11-09
\bibitem{NCKP99} J. Kaczorowski and A. Perelli. On the structure of the Selberg class, I: $0\le d\le1$. \emph{Acta Mathematica} 182 (1999), 207--241. % verified: 10.1007/bf02392574 2026-11-09
\bibitem{NCGrenanderSzego58} U. Grenander and G. Szeg\H o. \emph{Toeplitz Forms and Their Applications}. University of California Press, Berkeley, 1958. % verified: CrossRef 10.1525/9780520355408 2026-11-10
\end{thebibliography}
